% ============================================================
% Antecedentes de magnetopriming en vegetales
% Incluir en marco_teorico.tex con:
%   % ============================================================
% Antecedentes de magnetopriming en vegetales
% Incluir en marco_teorico.tex con:
%   % ============================================================
% Antecedentes de magnetopriming en vegetales
% Incluir en marco_teorico.tex con:
%   % ============================================================
% Antecedentes de magnetopriming en vegetales
% Incluir en marco_teorico.tex con:
%   \input{secciones/antecedentes}
% (o descomentar el include en tesis.tex si se prefiere como
%  sección separada)
% ============================================================

\subsection{Antecedentes de magnetopriming en vegetales}

\subsubsection{Visión general}

El magnetopriming --esto es, el tratamiento de semillas con campos magnéticos previo a la siembra-- ha recibido atención creciente en las últimas dos décadas como una técnica física de bajo costo y potencialmente reproducible para mejorar la germinación, el vigor y el rendimiento de diversos cultivos \cite{araujo_2016, sarraf_2020}. La evidencia acumulada abarca más de treinta especies vegetales, con resultados que van desde incrementos en la velocidad de germinación hasta mejoras en la tolerancia a estrés abiótico \cite{radhakrishnan_2019, farooq_2019}. Un meta-análisis reciente que integra 45 estudios y 29 especies confirma que los campos magnéticos no uniformes incrementan el peso fresco de las plántulas en un 123\,\% en promedio, mientras que los uniformes mejoran la germinación en un 12,8\,\% \cite{tapia_belmonte_2023}.

En el maíz (\textit{Zea mays}), la especie más estudiada en este contexto junto con la soja, existen al menos 18 estudios experimentales directos que documentan efectos del campo magnético estático (SMF) sobre la germinación, el crecimiento, la fisiología y el rendimiento \cite{aladjadjiyan_2002, florez_2007, vashisth_joshi_2017, ferroni_2024}. Los parámetros de exposición reportados abarcan intensidades desde 40 hasta 560\,mT y tiempos desde 1 minuto hasta exposición continua, con un consenso emergente en torno a los 200\,mT durante 1 hora como condición óptima \cite{shine_guruprasad_2012, vashisth_joshi_2017, kataria_2017}. Ferroni et al. \cite{ferroni_2024} observaron que 1 hora de exposición a 150\,mT aumentó la longitud total de las plántulas en un 108.9\% al décimo día, sin cambios significativos en la energía germinativa ni en el poder germinativo.

\subsubsection{Efectos sobre la germinación}

Cuando las semillas de maíz ya poseen alta viabilidad (>90\,\%), el SMF no siempre incrementa el porcentaje final de germinación, pero consistentemente \textbf{acelera el proceso}. El tiempo medio de germinación (MGT) se reduce entre un 8\,\% y un 34\,\% según la dosis y la variedad, y el inicio de la germinación (T1) se adelanta hasta 13 horas \cite{martinez_2017, florez_2007, torres_2019}. Aladjadjiyan \cite{aladjadjiyan_2002} reportó un incremento de la germinabilidad (GE) del 56\,\% al 80\,\% con 150\,mT durante 10 minutos, alcanzando un 100\,\% de germinación final frente al 85\,\% del control. Vashisth y Joshi \cite{vashisth_joshi_2017} encontraron que 200\,mT durante 1 hora elevó la germinación del 80\,\% al 96\,\% y aumentó la velocidad de germinación en un 28\,\%.

Torres et al. \cite{torres_2019} reportaron un efecto dosis-respuesta no lineal en forma de U invertida: la dosis más baja (50\,mT, 1 min) redujo el MGT un 12,4\,\%, mientras que dosis superiores a 150\,mT resultaron inhibitorias. Esto sugiere que no solo la intensidad, sino también la homogeneidad espacial del campo juega un rol determinante.

\subsubsection{Efectos sobre el crecimiento y el rendimiento}

Los incrementos en biomasa y rendimiento son posiblemente el resultado más impactante del magnetopriming en maíz. En condiciones de laboratorio, Shine y Guruprasad \cite{shine_guruprasad_2012} reportaron aumentos del 78\,\% en área foliar, 52\,\% en peso seco de brotes y 58\,\% en peso seco de raíces con 200\,mT durante 1 hora. Flórez et al. \cite{florez_2007} encontraron que exposiciones de 24 horas a 125--250\,mT incrementaron el peso fresco total un 34\,\% al décimo día.

En condiciones de campo, los efectos persisten al menos 30 días después de la siembra. Vashisth y Joshi \cite{vashisth_joshi_2017} reportaron incrementos del 21--23\,\% en el rendimiento de grano en parcelas experimentales, mientras que Siyami et al. \cite{siyami_2018} obtuvieron un 20\,\% de aumento bajo estrés hídrico moderado con exposiciones de solo 40\,mT durante 10 minutos. Particularmente notable es el desarrollo radical: Anand et al. \cite{anand_2012} observaron un incremento de hasta el 790\,\% en la longitud radicular total bajo estrés hídrico leve ($-0{,}2$\,MPa), sugiriendo que la mejora en la captación de agua es uno de los mecanismos principales.

\subsubsection{Efectos bioquímicos y fisiológicos}

El cuadro bioquímico que emerge de la literatura sigue un patrón temporal bifásico:

\begin{enumerate}
    \item \textbf{Fase temprana (0--72\,h de germinación):} El SMF induce un aumento controlado de especies reactivas de oxígeno (ROS). Shine et al. \cite{shine_2017} midieron incrementos del 34\,\% en O$_2^{\bullet-}$, 48\,\% en H$_2$O$_2$ y 39\,\% en $\cdot$OH en embriones de maíz expuestos a 200\,mT. La actividad de la enzima peroxidasa (POD) asociada a la pared celular aumentó un 58\,\%, facilitando la ruptura de la cubierta seminal y la elongación celular \cite{shine_2017, kataria_2017}.
    \item \textbf{Fase tardía (plántulas establecidas, 7--45 días):} Las plantas provenientes de semillas tratadas presentan \textbf{menores} niveles de ROS y menor actividad de enzimas antioxidantes (SOD, CAT, POX), indicando que el sistema antioxidante se ha vuelto más eficiente y la energía puede redirigirse al crecimiento \cite{shine_guruprasad_2012, anand_2012}.
\end{enumerate}

Simultáneamente, el SMF acelera la hidratación de la semilla. Vashisth y Nagarajan \cite{vashisth_nagarajan_2009} demostraron mediante isotermas de sorción que el tratamiento magnético incrementa los sitios de adsorción débil de agua hasta un 61\,\%, es decir, agua más disponible metabólicamente. Mediciones de resonancia magnética nuclear (NMR) confirmaron que el agua de hidratación de macromoléculas aparece 3 horas antes en semillas tratadas, y que la absorción total de agua en fase III de imbibición aumenta de 1,177 a 1,586\,g\,H$_2$O/g \cite{vashisth_nagarajan_2010}.

En el plano fotosintético, Shine y Guruprasad \cite{shine_guruprasad_2012} reportaron que el índice de rendimiento fotosintético (PIABS) se duplica (+103\,\%) en plantas tratadas, acompañado de un aumento del 49\,\% en la densidad de centros de reacción activos del PSII (RC/CSm) y del 16\,\% en clorofila total. Saletnik et al. \cite{saletnik_2022_static} reportaron incrementos superiores al 60\,\% en la tasa fotosintética neta en maíz.

\subsubsection{Efectos bajo estrés abiótico}

El magnetopriming ha demostrado ser particularmente efectivo cuando las semillas germinan bajo condiciones adversas. Bajo \textbf{estrés salino} (100\,mM NaCl), Kataria et al. \cite{kataria_2017} encontraron que el pre-tratamiento con 200\,mT durante 1 hora incrementó la germinación un 19\,\% y el vigor I un 41\,\%, con un aumento del 62\,\% en la absorción de agua. Bajo \textbf{estrés hídrico}, Anand et al. \cite{anand_2012} mostraron que las plántulas tratadas presentaban un potencial hídrico 21\,\% mayor, una tasa fotosintética neta superior y un 44--58\,\% menos de peroxidación lipídica. Siyami et al. \cite{siyami_2018} reportaron que el tratamiento magnético mantuvo un rendimiento de grano 13\,\% superior al control bajo condiciones de déficit hídrico. Bajo \textbf{estrés por frío}, Radhakrishnan \cite{radhakrishnan_2019} citó estudios donde campos de 150\,mT mejoraron la clorofila, los compuestos fenólicos y el intercambio gaseoso en maíz, reduciendo la permeabilidad de membranas.

\subsubsection{Mecanismos propuestos}

Si bien el mecanismo molecular completo no está dilucidado \cite{araujo_2016, maffei_2022}, se han propuesto varias vías que operan en distintas escalas:

\begin{enumerate}
    \item \textbf{Mecanismo de par radical (RPM) mediado por criptocromos:} Es el modelo más aceptado para la magnetorrecepción en plantas \cite{maffei_2022, saletnik_2022_static}. En los criptocromos (fotorreceptores de luz azul), la fotoexcitación del cofactor FAD genera un par radical cuyo espín electrónico es sensible al campo magnético. El par FADH$\bullet$/O$_2^{\bullet-}$ es el candidato propuesto en plantas. El RPM depende de la luz, pero también se observan efectos en oscuridad, sugiriendo vías independientes \cite{maffei_2022}.
    \item \textbf{Resonancia iónica de ciclotrón (ICR) del Ca$^{2+}$:} Propone que la frecuencia de Larmor del ion Ca$^{2+}$ puede coincidir con la frecuencia efectiva del campo, alterando el flujo iónico a través de membranas \cite{radhakrishnan_2019, saletnik_2022_static}. Belyavskaya \cite{belyavskaya_2004} demostró sobre-saturación de Ca$^{2+}$ citosólico en células expuestas a campos débiles.
    \item \textbf{Señalización por ROS/H$_2$O$_2$:} El SMF induce un aumento controlado de ROS que actúa como señal para activar la movilización de reservas, la ruptura de la cubierta seminal y la regulación del balance hormonal ABA/GA \cite{shine_2017, kataria_2017, sarraf_2020}.
    \item \textbf{Reorganización del agua celular:} El SMF modifica las propiedades de unión del agua, incrementando la fracción de agua débilmente unida y su movilidad molecular, lo que acelera la hidratación de macromoléculas y la activación enzimática \cite{vashisth_nagarajan_2009, vashisth_nagarajan_2010, torres_socorro_2018}.
    \item \textbf{Activación enzimática:} Incremento consistente de $\alpha$-amilasa (+17--22\,\%), proteasa (+5--8\,\%) y deshidrogenasa (+44--48\,\%) en semillas tratadas, proporcionando mayor energía y sustratos para el crecimiento temprano \cite{vashisth_nagarajan_2010, kataria_2017}.
\end{enumerate}

\subsubsection{Estado actual y brechas de conocimiento}

A pesar de los resultados prometedores, la literatura presenta limitaciones importantes \cite{araujo_2016, radhakrishnan_2019, sarraf_2020}: (i) falta de estandarización en los protocolos de exposición, (ii) pocos estudios que evalúen cambios en la expresión génica, (iii) escasa caracterización de la homogeneidad y el gradiente espacial del campo aplicado, y (iv) ausencia de estudios sobre persistencia transgeneracional de los efectos. Recientemente se ha señalado que el tipo de campo (uniforme vs.\ no uniforme) es quizás más determinante que la intensidad o la duración de la exposición \cite{tapia_belmonte_2023}, y que los campos con gradiente --como los generados por imanes toroidales-- producen sistemáticamente mejores resultados que los campos homogéneos \cite{torres_2019}.

\endinput

% (o descomentar el include en tesis.tex si se prefiere como
%  sección separada)
% ============================================================

\subsection{Antecedentes de magnetopriming en vegetales}

\subsubsection{Visión general}

El magnetopriming --esto es, el tratamiento de semillas con campos magnéticos previo a la siembra-- ha recibido atención creciente en las últimas dos décadas como una técnica física de bajo costo y potencialmente reproducible para mejorar la germinación, el vigor y el rendimiento de diversos cultivos \cite{araujo_2016, sarraf_2020}. La evidencia acumulada abarca más de treinta especies vegetales, con resultados que van desde incrementos en la velocidad de germinación hasta mejoras en la tolerancia a estrés abiótico \cite{radhakrishnan_2019, farooq_2019}. Un meta-análisis reciente que integra 45 estudios y 29 especies confirma que los campos magnéticos no uniformes incrementan el peso fresco de las plántulas en un 123\,\% en promedio, mientras que los uniformes mejoran la germinación en un 12,8\,\% \cite{tapia_belmonte_2023}.

En el maíz (\textit{Zea mays}), la especie más estudiada en este contexto junto con la soja, existen al menos 18 estudios experimentales directos que documentan efectos del campo magnético estático (SMF) sobre la germinación, el crecimiento, la fisiología y el rendimiento \cite{aladjadjiyan_2002, florez_2007, vashisth_joshi_2017, ferroni_2024}. Los parámetros de exposición reportados abarcan intensidades desde 40 hasta 560\,mT y tiempos desde 1 minuto hasta exposición continua, con un consenso emergente en torno a los 200\,mT durante 1 hora como condición óptima \cite{shine_guruprasad_2012, vashisth_joshi_2017, kataria_2017}. Ferroni et al. \cite{ferroni_2024} observaron que 1 hora de exposición a 150\,mT aumentó la longitud total de las plántulas en un 108.9\% al décimo día, sin cambios significativos en la energía germinativa ni en el poder germinativo.

\subsubsection{Efectos sobre la germinación}

Cuando las semillas de maíz ya poseen alta viabilidad (>90\,\%), el SMF no siempre incrementa el porcentaje final de germinación, pero consistentemente \textbf{acelera el proceso}. El tiempo medio de germinación (MGT) se reduce entre un 8\,\% y un 34\,\% según la dosis y la variedad, y el inicio de la germinación (T1) se adelanta hasta 13 horas \cite{martinez_2017, florez_2007, torres_2019}. Aladjadjiyan \cite{aladjadjiyan_2002} reportó un incremento de la germinabilidad (GE) del 56\,\% al 80\,\% con 150\,mT durante 10 minutos, alcanzando un 100\,\% de germinación final frente al 85\,\% del control. Vashisth y Joshi \cite{vashisth_joshi_2017} encontraron que 200\,mT durante 1 hora elevó la germinación del 80\,\% al 96\,\% y aumentó la velocidad de germinación en un 28\,\%.

Torres et al. \cite{torres_2019} reportaron un efecto dosis-respuesta no lineal en forma de U invertida: la dosis más baja (50\,mT, 1 min) redujo el MGT un 12,4\,\%, mientras que dosis superiores a 150\,mT resultaron inhibitorias. Esto sugiere que no solo la intensidad, sino también la homogeneidad espacial del campo juega un rol determinante.

\subsubsection{Efectos sobre el crecimiento y el rendimiento}

Los incrementos en biomasa y rendimiento son posiblemente el resultado más impactante del magnetopriming en maíz. En condiciones de laboratorio, Shine y Guruprasad \cite{shine_guruprasad_2012} reportaron aumentos del 78\,\% en área foliar, 52\,\% en peso seco de brotes y 58\,\% en peso seco de raíces con 200\,mT durante 1 hora. Flórez et al. \cite{florez_2007} encontraron que exposiciones de 24 horas a 125--250\,mT incrementaron el peso fresco total un 34\,\% al décimo día.

En condiciones de campo, los efectos persisten al menos 30 días después de la siembra. Vashisth y Joshi \cite{vashisth_joshi_2017} reportaron incrementos del 21--23\,\% en el rendimiento de grano en parcelas experimentales, mientras que Siyami et al. \cite{siyami_2018} obtuvieron un 20\,\% de aumento bajo estrés hídrico moderado con exposiciones de solo 40\,mT durante 10 minutos. Particularmente notable es el desarrollo radical: Anand et al. \cite{anand_2012} observaron un incremento de hasta el 790\,\% en la longitud radicular total bajo estrés hídrico leve ($-0{,}2$\,MPa), sugiriendo que la mejora en la captación de agua es uno de los mecanismos principales.

\subsubsection{Efectos bioquímicos y fisiológicos}

El cuadro bioquímico que emerge de la literatura sigue un patrón temporal bifásico:

\begin{enumerate}
    \item \textbf{Fase temprana (0--72\,h de germinación):} El SMF induce un aumento controlado de especies reactivas de oxígeno (ROS). Shine et al. \cite{shine_2017} midieron incrementos del 34\,\% en O$_2^{\bullet-}$, 48\,\% en H$_2$O$_2$ y 39\,\% en $\cdot$OH en embriones de maíz expuestos a 200\,mT. La actividad de la enzima peroxidasa (POD) asociada a la pared celular aumentó un 58\,\%, facilitando la ruptura de la cubierta seminal y la elongación celular \cite{shine_2017, kataria_2017}.
    \item \textbf{Fase tardía (plántulas establecidas, 7--45 días):} Las plantas provenientes de semillas tratadas presentan \textbf{menores} niveles de ROS y menor actividad de enzimas antioxidantes (SOD, CAT, POX), indicando que el sistema antioxidante se ha vuelto más eficiente y la energía puede redirigirse al crecimiento \cite{shine_guruprasad_2012, anand_2012}.
\end{enumerate}

Simultáneamente, el SMF acelera la hidratación de la semilla. Vashisth y Nagarajan \cite{vashisth_nagarajan_2009} demostraron mediante isotermas de sorción que el tratamiento magnético incrementa los sitios de adsorción débil de agua hasta un 61\,\%, es decir, agua más disponible metabólicamente. Mediciones de resonancia magnética nuclear (NMR) confirmaron que el agua de hidratación de macromoléculas aparece 3 horas antes en semillas tratadas, y que la absorción total de agua en fase III de imbibición aumenta de 1,177 a 1,586\,g\,H$_2$O/g \cite{vashisth_nagarajan_2010}.

En el plano fotosintético, Shine y Guruprasad \cite{shine_guruprasad_2012} reportaron que el índice de rendimiento fotosintético (PIABS) se duplica (+103\,\%) en plantas tratadas, acompañado de un aumento del 49\,\% en la densidad de centros de reacción activos del PSII (RC/CSm) y del 16\,\% en clorofila total. Saletnik et al. \cite{saletnik_2022_static} reportaron incrementos superiores al 60\,\% en la tasa fotosintética neta en maíz.

\subsubsection{Efectos bajo estrés abiótico}

El magnetopriming ha demostrado ser particularmente efectivo cuando las semillas germinan bajo condiciones adversas. Bajo \textbf{estrés salino} (100\,mM NaCl), Kataria et al. \cite{kataria_2017} encontraron que el pre-tratamiento con 200\,mT durante 1 hora incrementó la germinación un 19\,\% y el vigor I un 41\,\%, con un aumento del 62\,\% en la absorción de agua. Bajo \textbf{estrés hídrico}, Anand et al. \cite{anand_2012} mostraron que las plántulas tratadas presentaban un potencial hídrico 21\,\% mayor, una tasa fotosintética neta superior y un 44--58\,\% menos de peroxidación lipídica. Siyami et al. \cite{siyami_2018} reportaron que el tratamiento magnético mantuvo un rendimiento de grano 13\,\% superior al control bajo condiciones de déficit hídrico. Bajo \textbf{estrés por frío}, Radhakrishnan \cite{radhakrishnan_2019} citó estudios donde campos de 150\,mT mejoraron la clorofila, los compuestos fenólicos y el intercambio gaseoso en maíz, reduciendo la permeabilidad de membranas.

\subsubsection{Mecanismos propuestos}

Si bien el mecanismo molecular completo no está dilucidado \cite{araujo_2016, maffei_2022}, se han propuesto varias vías que operan en distintas escalas:

\begin{enumerate}
    \item \textbf{Mecanismo de par radical (RPM) mediado por criptocromos:} Es el modelo más aceptado para la magnetorrecepción en plantas \cite{maffei_2022, saletnik_2022_static}. En los criptocromos (fotorreceptores de luz azul), la fotoexcitación del cofactor FAD genera un par radical cuyo espín electrónico es sensible al campo magnético. El par FADH$\bullet$/O$_2^{\bullet-}$ es el candidato propuesto en plantas. El RPM depende de la luz, pero también se observan efectos en oscuridad, sugiriendo vías independientes \cite{maffei_2022}.
    \item \textbf{Resonancia iónica de ciclotrón (ICR) del Ca$^{2+}$:} Propone que la frecuencia de Larmor del ion Ca$^{2+}$ puede coincidir con la frecuencia efectiva del campo, alterando el flujo iónico a través de membranas \cite{radhakrishnan_2019, saletnik_2022_static}. Belyavskaya \cite{belyavskaya_2004} demostró sobre-saturación de Ca$^{2+}$ citosólico en células expuestas a campos débiles.
    \item \textbf{Señalización por ROS/H$_2$O$_2$:} El SMF induce un aumento controlado de ROS que actúa como señal para activar la movilización de reservas, la ruptura de la cubierta seminal y la regulación del balance hormonal ABA/GA \cite{shine_2017, kataria_2017, sarraf_2020}.
    \item \textbf{Reorganización del agua celular:} El SMF modifica las propiedades de unión del agua, incrementando la fracción de agua débilmente unida y su movilidad molecular, lo que acelera la hidratación de macromoléculas y la activación enzimática \cite{vashisth_nagarajan_2009, vashisth_nagarajan_2010, torres_socorro_2018}.
    \item \textbf{Activación enzimática:} Incremento consistente de $\alpha$-amilasa (+17--22\,\%), proteasa (+5--8\,\%) y deshidrogenasa (+44--48\,\%) en semillas tratadas, proporcionando mayor energía y sustratos para el crecimiento temprano \cite{vashisth_nagarajan_2010, kataria_2017}.
\end{enumerate}

\subsubsection{Estado actual y brechas de conocimiento}

A pesar de los resultados prometedores, la literatura presenta limitaciones importantes \cite{araujo_2016, radhakrishnan_2019, sarraf_2020}: (i) falta de estandarización en los protocolos de exposición, (ii) pocos estudios que evalúen cambios en la expresión génica, (iii) escasa caracterización de la homogeneidad y el gradiente espacial del campo aplicado, y (iv) ausencia de estudios sobre persistencia transgeneracional de los efectos. Recientemente se ha señalado que el tipo de campo (uniforme vs.\ no uniforme) es quizás más determinante que la intensidad o la duración de la exposición \cite{tapia_belmonte_2023}, y que los campos con gradiente --como los generados por imanes toroidales-- producen sistemáticamente mejores resultados que los campos homogéneos \cite{torres_2019}.

\endinput

% (o descomentar el include en tesis.tex si se prefiere como
%  sección separada)
% ============================================================

\subsection{Antecedentes de magnetopriming en vegetales}

\subsubsection{Visión general}

El magnetopriming --esto es, el tratamiento de semillas con campos magnéticos previo a la siembra-- ha recibido atención creciente en las últimas dos décadas como una técnica física de bajo costo y potencialmente reproducible para mejorar la germinación, el vigor y el rendimiento de diversos cultivos \cite{araujo_2016, sarraf_2020}. La evidencia acumulada abarca más de treinta especies vegetales, con resultados que van desde incrementos en la velocidad de germinación hasta mejoras en la tolerancia a estrés abiótico \cite{radhakrishnan_2019, farooq_2019}. Un meta-análisis reciente que integra 45 estudios y 29 especies confirma que los campos magnéticos no uniformes incrementan el peso fresco de las plántulas en un 123\,\% en promedio, mientras que los uniformes mejoran la germinación en un 12,8\,\% \cite{tapia_belmonte_2023}.

En el maíz (\textit{Zea mays}), la especie más estudiada en este contexto junto con la soja, existen al menos 18 estudios experimentales directos que documentan efectos del campo magnético estático (SMF) sobre la germinación, el crecimiento, la fisiología y el rendimiento \cite{aladjadjiyan_2002, florez_2007, vashisth_joshi_2017, ferroni_2024}. Los parámetros de exposición reportados abarcan intensidades desde 40 hasta 560\,mT y tiempos desde 1 minuto hasta exposición continua, con un consenso emergente en torno a los 200\,mT durante 1 hora como condición óptima \cite{shine_guruprasad_2012, vashisth_joshi_2017, kataria_2017}. Ferroni et al. \cite{ferroni_2024} observaron que 1 hora de exposición a 150\,mT aumentó la longitud total de las plántulas en un 108.9\% al décimo día, sin cambios significativos en la energía germinativa ni en el poder germinativo.

\subsubsection{Efectos sobre la germinación}

Cuando las semillas de maíz ya poseen alta viabilidad (>90\,\%), el SMF no siempre incrementa el porcentaje final de germinación, pero consistentemente \textbf{acelera el proceso}. El tiempo medio de germinación (MGT) se reduce entre un 8\,\% y un 34\,\% según la dosis y la variedad, y el inicio de la germinación (T1) se adelanta hasta 13 horas \cite{martinez_2017, florez_2007, torres_2019}. Aladjadjiyan \cite{aladjadjiyan_2002} reportó un incremento de la germinabilidad (GE) del 56\,\% al 80\,\% con 150\,mT durante 10 minutos, alcanzando un 100\,\% de germinación final frente al 85\,\% del control. Vashisth y Joshi \cite{vashisth_joshi_2017} encontraron que 200\,mT durante 1 hora elevó la germinación del 80\,\% al 96\,\% y aumentó la velocidad de germinación en un 28\,\%.

Torres et al. \cite{torres_2019} reportaron un efecto dosis-respuesta no lineal en forma de U invertida: la dosis más baja (50\,mT, 1 min) redujo el MGT un 12,4\,\%, mientras que dosis superiores a 150\,mT resultaron inhibitorias. Esto sugiere que no solo la intensidad, sino también la homogeneidad espacial del campo juega un rol determinante.

\subsubsection{Efectos sobre el crecimiento y el rendimiento}

Los incrementos en biomasa y rendimiento son posiblemente el resultado más impactante del magnetopriming en maíz. En condiciones de laboratorio, Shine y Guruprasad \cite{shine_guruprasad_2012} reportaron aumentos del 78\,\% en área foliar, 52\,\% en peso seco de brotes y 58\,\% en peso seco de raíces con 200\,mT durante 1 hora. Flórez et al. \cite{florez_2007} encontraron que exposiciones de 24 horas a 125--250\,mT incrementaron el peso fresco total un 34\,\% al décimo día.

En condiciones de campo, los efectos persisten al menos 30 días después de la siembra. Vashisth y Joshi \cite{vashisth_joshi_2017} reportaron incrementos del 21--23\,\% en el rendimiento de grano en parcelas experimentales, mientras que Siyami et al. \cite{siyami_2018} obtuvieron un 20\,\% de aumento bajo estrés hídrico moderado con exposiciones de solo 40\,mT durante 10 minutos. Particularmente notable es el desarrollo radical: Anand et al. \cite{anand_2012} observaron un incremento de hasta el 790\,\% en la longitud radicular total bajo estrés hídrico leve ($-0{,}2$\,MPa), sugiriendo que la mejora en la captación de agua es uno de los mecanismos principales.

\subsubsection{Efectos bioquímicos y fisiológicos}

El cuadro bioquímico que emerge de la literatura sigue un patrón temporal bifásico:

\begin{enumerate}
    \item \textbf{Fase temprana (0--72\,h de germinación):} El SMF induce un aumento controlado de especies reactivas de oxígeno (ROS). Shine et al. \cite{shine_2017} midieron incrementos del 34\,\% en O$_2^{\bullet-}$, 48\,\% en H$_2$O$_2$ y 39\,\% en $\cdot$OH en embriones de maíz expuestos a 200\,mT. La actividad de la enzima peroxidasa (POD) asociada a la pared celular aumentó un 58\,\%, facilitando la ruptura de la cubierta seminal y la elongación celular \cite{shine_2017, kataria_2017}.
    \item \textbf{Fase tardía (plántulas establecidas, 7--45 días):} Las plantas provenientes de semillas tratadas presentan \textbf{menores} niveles de ROS y menor actividad de enzimas antioxidantes (SOD, CAT, POX), indicando que el sistema antioxidante se ha vuelto más eficiente y la energía puede redirigirse al crecimiento \cite{shine_guruprasad_2012, anand_2012}.
\end{enumerate}

Simultáneamente, el SMF acelera la hidratación de la semilla. Vashisth y Nagarajan \cite{vashisth_nagarajan_2009} demostraron mediante isotermas de sorción que el tratamiento magnético incrementa los sitios de adsorción débil de agua hasta un 61\,\%, es decir, agua más disponible metabólicamente. Mediciones de resonancia magnética nuclear (NMR) confirmaron que el agua de hidratación de macromoléculas aparece 3 horas antes en semillas tratadas, y que la absorción total de agua en fase III de imbibición aumenta de 1,177 a 1,586\,g\,H$_2$O/g \cite{vashisth_nagarajan_2010}.

En el plano fotosintético, Shine y Guruprasad \cite{shine_guruprasad_2012} reportaron que el índice de rendimiento fotosintético (PIABS) se duplica (+103\,\%) en plantas tratadas, acompañado de un aumento del 49\,\% en la densidad de centros de reacción activos del PSII (RC/CSm) y del 16\,\% en clorofila total. Saletnik et al. \cite{saletnik_2022_static} reportaron incrementos superiores al 60\,\% en la tasa fotosintética neta en maíz.

\subsubsection{Efectos bajo estrés abiótico}

El magnetopriming ha demostrado ser particularmente efectivo cuando las semillas germinan bajo condiciones adversas. Bajo \textbf{estrés salino} (100\,mM NaCl), Kataria et al. \cite{kataria_2017} encontraron que el pre-tratamiento con 200\,mT durante 1 hora incrementó la germinación un 19\,\% y el vigor I un 41\,\%, con un aumento del 62\,\% en la absorción de agua. Bajo \textbf{estrés hídrico}, Anand et al. \cite{anand_2012} mostraron que las plántulas tratadas presentaban un potencial hídrico 21\,\% mayor, una tasa fotosintética neta superior y un 44--58\,\% menos de peroxidación lipídica. Siyami et al. \cite{siyami_2018} reportaron que el tratamiento magnético mantuvo un rendimiento de grano 13\,\% superior al control bajo condiciones de déficit hídrico. Bajo \textbf{estrés por frío}, Radhakrishnan \cite{radhakrishnan_2019} citó estudios donde campos de 150\,mT mejoraron la clorofila, los compuestos fenólicos y el intercambio gaseoso en maíz, reduciendo la permeabilidad de membranas.

\subsubsection{Mecanismos propuestos}

Si bien el mecanismo molecular completo no está dilucidado \cite{araujo_2016, maffei_2022}, se han propuesto varias vías que operan en distintas escalas:

\begin{enumerate}
    \item \textbf{Mecanismo de par radical (RPM) mediado por criptocromos:} Es el modelo más aceptado para la magnetorrecepción en plantas \cite{maffei_2022, saletnik_2022_static}. En los criptocromos (fotorreceptores de luz azul), la fotoexcitación del cofactor FAD genera un par radical cuyo espín electrónico es sensible al campo magnético. El par FADH$\bullet$/O$_2^{\bullet-}$ es el candidato propuesto en plantas. El RPM depende de la luz, pero también se observan efectos en oscuridad, sugiriendo vías independientes \cite{maffei_2022}.
    \item \textbf{Resonancia iónica de ciclotrón (ICR) del Ca$^{2+}$:} Propone que la frecuencia de Larmor del ion Ca$^{2+}$ puede coincidir con la frecuencia efectiva del campo, alterando el flujo iónico a través de membranas \cite{radhakrishnan_2019, saletnik_2022_static}. Belyavskaya \cite{belyavskaya_2004} demostró sobre-saturación de Ca$^{2+}$ citosólico en células expuestas a campos débiles.
    \item \textbf{Señalización por ROS/H$_2$O$_2$:} El SMF induce un aumento controlado de ROS que actúa como señal para activar la movilización de reservas, la ruptura de la cubierta seminal y la regulación del balance hormonal ABA/GA \cite{shine_2017, kataria_2017, sarraf_2020}.
    \item \textbf{Reorganización del agua celular:} El SMF modifica las propiedades de unión del agua, incrementando la fracción de agua débilmente unida y su movilidad molecular, lo que acelera la hidratación de macromoléculas y la activación enzimática \cite{vashisth_nagarajan_2009, vashisth_nagarajan_2010, torres_socorro_2018}.
    \item \textbf{Activación enzimática:} Incremento consistente de $\alpha$-amilasa (+17--22\,\%), proteasa (+5--8\,\%) y deshidrogenasa (+44--48\,\%) en semillas tratadas, proporcionando mayor energía y sustratos para el crecimiento temprano \cite{vashisth_nagarajan_2010, kataria_2017}.
\end{enumerate}

\subsubsection{Estado actual y brechas de conocimiento}

A pesar de los resultados prometedores, la literatura presenta limitaciones importantes \cite{araujo_2016, radhakrishnan_2019, sarraf_2020}: (i) falta de estandarización en los protocolos de exposición, (ii) pocos estudios que evalúen cambios en la expresión génica, (iii) escasa caracterización de la homogeneidad y el gradiente espacial del campo aplicado, y (iv) ausencia de estudios sobre persistencia transgeneracional de los efectos. Recientemente se ha señalado que el tipo de campo (uniforme vs.\ no uniforme) es quizás más determinante que la intensidad o la duración de la exposición \cite{tapia_belmonte_2023}, y que los campos con gradiente --como los generados por imanes toroidales-- producen sistemáticamente mejores resultados que los campos homogéneos \cite{torres_2019}.

\endinput

% (o descomentar el include en tesis.tex si se prefiere como
%  sección separada)
% ============================================================

\subsection{Antecedentes de magnetopriming en vegetales}

\subsubsection{Visión general}

El magnetopriming --esto es, el tratamiento de semillas con campos magnéticos previo a la siembra-- ha recibido atención creciente en las últimas dos décadas como una técnica física de bajo costo y potencialmente reproducible para mejorar la germinación, el vigor y el rendimiento de diversos cultivos \cite{araujo_2016, sarraf_2020}. La evidencia acumulada abarca más de treinta especies vegetales, con resultados que van desde incrementos en la velocidad de germinación hasta mejoras en la tolerancia a estrés abiótico \cite{radhakrishnan_2019, farooq_2019}. Un meta-análisis reciente que integra 45 estudios y 29 especies confirma que los campos magnéticos no uniformes incrementan el peso fresco de las plántulas en un 123\,\% en promedio, mientras que los uniformes mejoran la germinación en un 12,8\,\% \cite{tapia_belmonte_2023}.

En el maíz (\textit{Zea mays}), la especie más estudiada en este contexto junto con la soja, existen al menos 18 estudios experimentales directos que documentan efectos del campo magnético estático (SMF) sobre la germinación, el crecimiento, la fisiología y el rendimiento \cite{aladjadjiyan_2002, florez_2007, vashisth_joshi_2017, ferroni_2024}. Los parámetros de exposición reportados abarcan intensidades desde 40 hasta 560\,mT y tiempos desde 1 minuto hasta exposición continua, con un consenso emergente en torno a los 200\,mT durante 1 hora como condición óptima \cite{shine_guruprasad_2012, vashisth_joshi_2017, kataria_2017}. Ferroni et al. \cite{ferroni_2024} observaron que 1 hora de exposición a 150\,mT aumentó la longitud total de las plántulas en un 108.9\% al décimo día, sin cambios significativos en la energía germinativa ni en el poder germinativo.

\subsubsection{Efectos sobre la germinación}

Cuando las semillas de maíz ya poseen alta viabilidad (>90\,\%), el SMF no siempre incrementa el porcentaje final de germinación, pero consistentemente \textbf{acelera el proceso}. El tiempo medio de germinación (MGT) se reduce entre un 8\,\% y un 34\,\% según la dosis y la variedad, y el inicio de la germinación (T1) se adelanta hasta 13 horas \cite{martinez_2017, florez_2007, torres_2019}. Aladjadjiyan \cite{aladjadjiyan_2002} reportó un incremento de la germinabilidad (GE) del 56\,\% al 80\,\% con 150\,mT durante 10 minutos, alcanzando un 100\,\% de germinación final frente al 85\,\% del control. Vashisth y Joshi \cite{vashisth_joshi_2017} encontraron que 200\,mT durante 1 hora elevó la germinación del 80\,\% al 96\,\% y aumentó la velocidad de germinación en un 28\,\%.

Torres et al. \cite{torres_2019} reportaron un efecto dosis-respuesta no lineal en forma de U invertida: la dosis más baja (50\,mT, 1 min) redujo el MGT un 12,4\,\%, mientras que dosis superiores a 150\,mT resultaron inhibitorias. Esto sugiere que no solo la intensidad, sino también la homogeneidad espacial del campo juega un rol determinante.

\subsubsection{Efectos sobre el crecimiento y el rendimiento}

Los incrementos en biomasa y rendimiento son posiblemente el resultado más impactante del magnetopriming en maíz. En condiciones de laboratorio, Shine y Guruprasad \cite{shine_guruprasad_2012} reportaron aumentos del 78\,\% en área foliar, 52\,\% en peso seco de brotes y 58\,\% en peso seco de raíces con 200\,mT durante 1 hora. Flórez et al. \cite{florez_2007} encontraron que exposiciones de 24 horas a 125--250\,mT incrementaron el peso fresco total un 34\,\% al décimo día.

En condiciones de campo, los efectos persisten al menos 30 días después de la siembra. Vashisth y Joshi \cite{vashisth_joshi_2017} reportaron incrementos del 21--23\,\% en el rendimiento de grano en parcelas experimentales, mientras que Siyami et al. \cite{siyami_2018} obtuvieron un 20\,\% de aumento bajo estrés hídrico moderado con exposiciones de solo 40\,mT durante 10 minutos. Particularmente notable es el desarrollo radical: Anand et al. \cite{anand_2012} observaron un incremento de hasta el 790\,\% en la longitud radicular total bajo estrés hídrico leve ($-0{,}2$\,MPa), sugiriendo que la mejora en la captación de agua es uno de los mecanismos principales.

\subsubsection{Efectos bioquímicos y fisiológicos}

El cuadro bioquímico que emerge de la literatura sigue un patrón temporal bifásico:

\begin{enumerate}
    \item \textbf{Fase temprana (0--72\,h de germinación):} El SMF induce un aumento controlado de especies reactivas de oxígeno (ROS). Shine et al. \cite{shine_2017} midieron incrementos del 34\,\% en O$_2^{\bullet-}$, 48\,\% en H$_2$O$_2$ y 39\,\% en $\cdot$OH en embriones de maíz expuestos a 200\,mT. La actividad de la enzima peroxidasa (POD) asociada a la pared celular aumentó un 58\,\%, facilitando la ruptura de la cubierta seminal y la elongación celular \cite{shine_2017, kataria_2017}.
    \item \textbf{Fase tardía (plántulas establecidas, 7--45 días):} Las plantas provenientes de semillas tratadas presentan \textbf{menores} niveles de ROS y menor actividad de enzimas antioxidantes (SOD, CAT, POX), indicando que el sistema antioxidante se ha vuelto más eficiente y la energía puede redirigirse al crecimiento \cite{shine_guruprasad_2012, anand_2012}.
\end{enumerate}

Simultáneamente, el SMF acelera la hidratación de la semilla. Vashisth y Nagarajan \cite{vashisth_nagarajan_2009} demostraron mediante isotermas de sorción que el tratamiento magnético incrementa los sitios de adsorción débil de agua hasta un 61\,\%, es decir, agua más disponible metabólicamente. Mediciones de resonancia magnética nuclear (NMR) confirmaron que el agua de hidratación de macromoléculas aparece 3 horas antes en semillas tratadas, y que la absorción total de agua en fase III de imbibición aumenta de 1,177 a 1,586\,g\,H$_2$O/g \cite{vashisth_nagarajan_2010}.

En el plano fotosintético, Shine y Guruprasad \cite{shine_guruprasad_2012} reportaron que el índice de rendimiento fotosintético (PIABS) se duplica (+103\,\%) en plantas tratadas, acompañado de un aumento del 49\,\% en la densidad de centros de reacción activos del PSII (RC/CSm) y del 16\,\% en clorofila total. Saletnik et al. \cite{saletnik_2022_static} reportaron incrementos superiores al 60\,\% en la tasa fotosintética neta en maíz.

\subsubsection{Efectos bajo estrés abiótico}

El magnetopriming ha demostrado ser particularmente efectivo cuando las semillas germinan bajo condiciones adversas. Bajo \textbf{estrés salino} (100\,mM NaCl), Kataria et al. \cite{kataria_2017} encontraron que el pre-tratamiento con 200\,mT durante 1 hora incrementó la germinación un 19\,\% y el vigor I un 41\,\%, con un aumento del 62\,\% en la absorción de agua. Bajo \textbf{estrés hídrico}, Anand et al. \cite{anand_2012} mostraron que las plántulas tratadas presentaban un potencial hídrico 21\,\% mayor, una tasa fotosintética neta superior y un 44--58\,\% menos de peroxidación lipídica. Siyami et al. \cite{siyami_2018} reportaron que el tratamiento magnético mantuvo un rendimiento de grano 13\,\% superior al control bajo condiciones de déficit hídrico. Bajo \textbf{estrés por frío}, Radhakrishnan \cite{radhakrishnan_2019} citó estudios donde campos de 150\,mT mejoraron la clorofila, los compuestos fenólicos y el intercambio gaseoso en maíz, reduciendo la permeabilidad de membranas.

\subsubsection{Mecanismos propuestos}

Si bien el mecanismo molecular completo no está dilucidado \cite{araujo_2016, maffei_2022}, se han propuesto varias vías que operan en distintas escalas:

\begin{enumerate}
    \item \textbf{Mecanismo de par radical (RPM) mediado por criptocromos:} Es el modelo más aceptado para la magnetorrecepción en plantas \cite{maffei_2022, saletnik_2022_static}. En los criptocromos (fotorreceptores de luz azul), la fotoexcitación del cofactor FAD genera un par radical cuyo espín electrónico es sensible al campo magnético. El par FADH$\bullet$/O$_2^{\bullet-}$ es el candidato propuesto en plantas. El RPM depende de la luz, pero también se observan efectos en oscuridad, sugiriendo vías independientes \cite{maffei_2022}.
    \item \textbf{Resonancia iónica de ciclotrón (ICR) del Ca$^{2+}$:} Propone que la frecuencia de Larmor del ion Ca$^{2+}$ puede coincidir con la frecuencia efectiva del campo, alterando el flujo iónico a través de membranas \cite{radhakrishnan_2019, saletnik_2022_static}. Belyavskaya \cite{belyavskaya_2004} demostró sobre-saturación de Ca$^{2+}$ citosólico en células expuestas a campos débiles.
    \item \textbf{Señalización por ROS/H$_2$O$_2$:} El SMF induce un aumento controlado de ROS que actúa como señal para activar la movilización de reservas, la ruptura de la cubierta seminal y la regulación del balance hormonal ABA/GA \cite{shine_2017, kataria_2017, sarraf_2020}.
    \item \textbf{Reorganización del agua celular:} El SMF modifica las propiedades de unión del agua, incrementando la fracción de agua débilmente unida y su movilidad molecular, lo que acelera la hidratación de macromoléculas y la activación enzimática \cite{vashisth_nagarajan_2009, vashisth_nagarajan_2010, torres_socorro_2018}.
    \item \textbf{Activación enzimática:} Incremento consistente de $\alpha$-amilasa (+17--22\,\%), proteasa (+5--8\,\%) y deshidrogenasa (+44--48\,\%) en semillas tratadas, proporcionando mayor energía y sustratos para el crecimiento temprano \cite{vashisth_nagarajan_2010, kataria_2017}.
\end{enumerate}

\subsubsection{Estado actual y brechas de conocimiento}

A pesar de los resultados prometedores, la literatura presenta limitaciones importantes \cite{araujo_2016, radhakrishnan_2019, sarraf_2020}: (i) falta de estandarización en los protocolos de exposición, (ii) pocos estudios que evalúen cambios en la expresión génica, (iii) escasa caracterización de la homogeneidad y el gradiente espacial del campo aplicado, y (iv) ausencia de estudios sobre persistencia transgeneracional de los efectos. Recientemente se ha señalado que el tipo de campo (uniforme vs.\ no uniforme) es quizás más determinante que la intensidad o la duración de la exposición \cite{tapia_belmonte_2023}, y que los campos con gradiente --como los generados por imanes toroidales-- producen sistemáticamente mejores resultados que los campos homogéneos \cite{torres_2019}.

\endinput
